Thermal stability is the property that allows a nuclear reactor to respond to
perturbations in power, temperature, and reactivity without diverging into an
uncontrolled excursion. It is the single most important dynamic characteristic
distinguishing a reactor that is safe by its very physics from one that depends
on engineered intervention to survive a transient. This monograph develops the
theory of reactor thermal stability from first principles and then uses it to
dissect the most consequential reactor accident in history---the destruction of
Chernobyl Unit~4 on 26~April 1986.

The treatment begins with the neutron balance and the point reactor kinetics
equations, including the central role of delayed neutrons, the inhour equation,
and the prompt-jump approximation. It then introduces the thermal source term,
decay heat, and the conduction--convection chain that carries fission energy
from fuel to coolant. With this groundwork, the book defines the reactivity
feedback coefficients---Doppler (fuel temperature), moderator and coolant
temperature, and the void coefficient---and shows how their signs and magnitudes
determine whether the coupled neutronic--thermal-hydraulic system is
self-correcting or self-amplifying.

A dedicated treatment of stability theory follows: linearisation of the coupled
equations, transfer functions and the Nyquist criterion, the nonlinear
Nordheim--Fuchs model of a self-limiting power excursion, and xenon-induced
spatial oscillations. The historical part surveys earlier stability-related
events (SL-1, Fermi-1, BWR oscillation events) before turning to the RBMK-1000
reactor, whose positive void coefficient and flawed control-rod design set the
stage for Chernobyl. The accident is reconstructed step by step and then
analysed through the lens of the feedback theory developed earlier, including
the evolution of the official understanding from INSAG-1 (1986) to INSAG-7
(1992). The work closes with the principles of inherent and passive safety in
modern reactor design and a computational demonstration using a companion
open-source point-kinetics toolkit.

Throughout, the emphasis is on the physical meaning of the mathematics:
\emph{a reactor is thermally stable when rising temperature removes reactivity
faster than rising power can add it.} Chernobyl is, in this language, the story
of a machine in which that inequality was reversed.
